\begin{abstract}
Distributed archaeological landscapes require spatially explicit monitoring because limited field capacity must be directed among components exposed to different combinations of land-surface change, terrain, drainage and climate forcing. This problem is acute at the Taxila World Heritage property in Pakistan, where multiple archaeological components occur within a heterogeneous peri-urban landscape. Existing geospatial assessments rarely combine comparable multi-temporal observations with acquisition-window weather, scale sensitivity, indicator redundancy and uncertainty in the resulting inspection order. This study developed the Climate-linked Heritage Inspection Prioritisation (CHIP) framework to integrate those evidence streams while preserving the distinction between landscape pressure and verified heritage condition. The analysis used 55 Landsat scenes organised into five post-monsoon epochs from 2003--2005 to 2023--2025, daily gridded weather for 1991--2025, approximately 30~m elevation data, and WorldCover-derived proxy land-cover labels. Seventeen of 18 inventoried components had usable geometry. Satellite indicators were compared on common support, climate trends and antecedent acquisition windows were characterised, terrain and drainage features were aggregated at 250, 500 and 1,000~m, and a hierarchical relative-priority score was tested through ablation, 50,000 weight perturbations and 2,000 spatial-block resamples. Common endpoint support comprised 386,280 cells, or 95.7371\% of the fixed grid; 16.9284\% met at least two adverse spectral criteria and 1.4220\% met all three. At the primary 500~m support, the Giri complex had the highest fixed score, 0.591176, but its central 95\% rank interval under weight perturbation extended from 1 to 9. Spatial resampling instead gave Bhallar median rank 1 and Giri median rank 2, with overlapping intervals. The 500 and 1,000~m rankings agreed more closely (Spearman $\rho=0.91846$) than the 250 and 1,000~m rankings ($\rho=0.71078$). A geographically buffered proxy experiment achieved a macro-averaged F1 score of 0.830795 on 4,188 held-out samples, but measured agreement with proxy land-cover labels only. CHIP therefore supports tiered field-inspection sequencing and mechanism-specific survey questions rather than an unqualified single-site ranking. It does not demonstrate monument damage, deterioration probability or a legally defined risk zone.
\end{abstract}
\noindent\textbf{Keywords:} cultural heritage; archaeological landscapes; Earth observation; multi-temporal Landsat; climate variability; terrain hydrology; spatial uncertainty; sensitivity analysis; field-inspection prioritisation; Taxila.
